\documentclass[12pt]{article}
\usepackage[utf8]{inputenc}
\usepackage[T1]{fontenc}
\usepackage{hyperref}
\title{Teor Language Resources for K-12 Curriculum Development}
\author{Compiled by Claude Sonnet 5 (claude-sonnet-5) for Andr\'as Kornai}
\date{}
\begin{document}
\maketitle

\section{Introduction}
Teor (ISO 639-3 tev) is an Austronesian language of the Watubela archipelago, Maluku, Indonesia, and is usually grouped with Kur as ``Teor-Kur''. Speaker counts in the sources vary from about 1,000 to about 4,500 and are dated. One source gives the Glottocode as teor1239, whereas the survey prompt gives teor1240. This should be checked. The search found only encyclopedic and survey pages.

\section{Primary Educational Resources}
None found.

\section{Linguistic and Reference Materials}
\begin{itemize}
\item Teor-Kur language, Wikipedia: \url{https://en.wikipedia.org/wiki/Teor-Kur_language}
\item Teor (island), Wikipedia: \url{https://en.wikipedia.org/wiki/Teor_(island)}
\item Teor numerals page, Max Planck Institute for Evolutionary Anthropology: \url{https://lingweb.eva.mpg.de/channumerals/Teor.htm}
\end{itemize}

\section{Religious and Community Resources}
\begin{itemize}
\item Joshua Project profile of the Teor in Indonesia: \url{https://joshuaproject.net/people_groups/15345}
\end{itemize}
This is a demographic profile only.

\section{Digital and Audio Resources}
None found.

\section{Curriculum Development Context}
No dictionary, grammar, primer, scripture or school programme for Teor turned up. Any curriculum would require new documentation work with the community. Kur is reportedly mutually intelligible with Teor and may be a useful partner variety.
\end{document}